\documentclass[../../main.tex]{subfiles}
\graphicspath{{../../image/}{../../figure/}}
\begin{document}

\section{可容许 $BV$ 解与广义特征线}\label{section:ch11-可容许BV解与广义特征线}

本章接下来的八节将讨论标量守恒律
\begin{align}
\partial_t u(x,t)+\partial_x f(u(x,t))=0,\label{align:ch11-11-1-1}
\end{align}
它是真正非线性的, 即 $f''(u)>0$, $-\infty<u<\infty$. 全文假设 $u$ 是 $(-\infty,\infty)\times[0,\infty)$ 上的可容许弱解, 其初值 $u_0$ 有界且在 $(-\infty,\infty)$ 上局部有界变差。由\cref{chapter:ch06-标量守恒律L1理论}可知 $u\in BV_{\mathrm{loc}}$, 且对任意 $t\in[0,\infty)$, 函数 $u(\cdot,t)$ 在 $(-\infty,\infty)$ 上局部有界变差。


如\cref{chapter:ch08-容许激波}所述, 对任意凸熵-熵流对 $(\eta,q)$, 熵激波可容许性条件将在解 $u$ 的激波集 $\mathscr{J}$ 上关于一维 Hausdorff 测度几乎处处满足。这进而推出 Lax $E$-条件也在 $\mathscr{J}$ 上几乎处处成立。因此, 我们有

\begin{align}
u(x+,t)\leqslant u(x-,t),\label{align:ch11-11-1-2}
\end{align}
对几乎所有 $t\in(0,\infty)$ 和所有 $x\in(-\infty,\infty)$ 成立。

由\cref{chapter:ch10-广义特征线}可知, 一条定义在时间区间 $[\sigma,\tau]\subset[0,\infty)$ 上的 Lipschitz 曲线 $\xi(\cdot)$, 若对几乎所有 $t\in[\sigma,\tau]$ 满足

\begin{align}
\dot{\xi}(t)=
\begin{cases}
f'(u(\xi(t)\pm,t)), & \text{当 }u(\xi(t)+,t)=u(\xi(t)-,t),\\[2mm]
\dfrac{f(u(\xi(t)+,t))-f(u(\xi(t)-,t))}{u(\xi(t)+,t)-u(\xi(t)-,t)}, & \text{当 }u(\xi(t)+,t)<u(\xi(t)-,t),
\end{cases}\label{align:ch11-11-1-3}
\end{align}
则称 $\xi(\cdot)$ 为与解 $u$ 关联的一条广义特征线。

真正非线性标量守恒律的特殊之处在于: 无激波的广义特征线本质上是经典特征线。

\begin{theorem}\label{theorem:ch11-11-1}
设 $\xi(\cdot)$ 是\eqref{align:ch11-11-1-1}的一条广义特征线, 与可容许解 $u$ 关联, 在时间区间 $[\sigma,\tau]$ 上无激波。则存在常数 $\bar u$ 使得
\begin{align}
u(\xi(\tau)+,\tau)\leqslant\bar u\leqslant u(\xi(\tau)-,\tau),\label{align:ch11-11-1-4}\\
u(\xi(t)+,t)=\bar u=u(\xi(t)-,t),\quad \sigma<t<\tau,\label{align:ch11-11-1-5}\\
u(\xi(\sigma)-,\sigma)\leqslant\bar u\leqslant u(\xi(\sigma)+,\sigma).\label{align:ch11-11-1-6}
\end{align}
特别地, $\xi(\cdot)$ 的图像是一条斜率为 $f'(\bar u)$ 的直线段。
\end{theorem}

\begin{proof}
固定 $r$ 和 $s$, $\sigma\leqslant r<s\leqslant\tau$。对 $\varepsilon>0$, 守恒律在区域 $\{(x,t):r<t<s,\xi(t)-\varepsilon<x<\xi(t)\}$ 和 $\{(x,t):r<t<s,\xi(t)<x<\xi(t)+\varepsilon\}$ 上分别取如下形式:
\begin{align}
&\int_{\xi(s)-\varepsilon}^{\xi(s)}u(x,s)\,\mathrm{d}x-\int_{\xi(r)-\varepsilon}^{\xi(r)}u(x,r)\,\mathrm{d}x\nonumber\\
&\quad=\int_r^s\{f(u(\xi(t)-\varepsilon+,t))-f(u(\xi(t)-,t))-\dot{\xi}(t)[u(\xi(t)-\varepsilon+,t)-u(\xi(t)-,t)]\}\,\mathrm{d}t,\label{align:ch11-11-1-7}
\end{align}
\begin{align}
&\int_{\xi(r)}^{\xi(r)+\varepsilon}u(x,r)\,\mathrm{d}x-\int_{\xi(s)}^{\xi(s)+\varepsilon}u(x,s)\,\mathrm{d}x\nonumber\\
&\quad=\int_r^s\{f(u(\xi(t)+\varepsilon-,t))-f(u(\xi(t)+,t))-\dot{\xi}(t)[u(\xi(t)+\varepsilon-,t)-u(\xi(t)+,t)]\}\,\mathrm{d}t.\label{align:ch11-11-1-8}
\end{align}
由\cref{chapter:ch10-广义特征线}可知, $\dot{\xi}(t)=f'(u(\xi(t)\pm,t))$ a.e. 于 $[r,s]$。于是由 $f$ 的凸性, \eqref{align:ch11-11-1-7} 与 \eqref{align:ch11-11-1-8} 的右端均非负。因此, 将两式乘以 $1/\varepsilon$ 并令 $\varepsilon\downarrow0$ 得到

\begin{align}
u(\xi(s)-,s)\geqslant u(\xi(r)-,r),\quad u(\xi(s)+,s)\leqslant u(\xi(r)+,r),\quad \sigma\leqslant r<s\leqslant\tau.\label{align:ch11-11-1-9}
\end{align}

现在固定 $t_1$ 和 $t_2$, $\sigma<t_1<t_2<\tau$, 使得 $u(\xi(t_1)-,t_1)=u(\xi(t_1)+,t_1)$, $u(\xi(t_2)-,t_2)=u(\xi(t_2)+,t_2)$。对任意固定的 $t\in(t_1,t_2)$, 先取 $r=t_1,s=t_2$, 再取 $r=t_1,s=t$, 最后取 $r=t,s=t_2$, 应用 \eqref{align:ch11-11-1-9}。由此得到 \eqref{align:ch11-11-1-5}。为完成证明, 对 $s=\tau$, $r\in(\sigma,\tau)$ 应用 \eqref{align:ch11-11-1-9} 得 \eqref{align:ch11-11-1-4}; 对 $r=\sigma$, $s\in(\sigma,\tau)$ 应用则得 \eqref{align:ch11-11-1-6}。
\end{proof}

\begin{corollary}\label{corollary:ch11-11-2}
假设 $\xi(\cdot)$ 与 $\zeta(\cdot)$ 是\eqref{align:ch11-11-1-1}的两条不同的广义特征线, 与可容许弱解 $u$ 关联, 且在时间区间 $[\sigma,\tau]$ 上无激波。则对任意 $t\in(\sigma,\tau)$, $\xi(\cdot)$ 与 $\zeta(\cdot)$ 不相交。
\end{corollary}

上述两个命题对极值后向特征线有重要推论。

\begin{theorem}\label{theorem:ch11-11-3}
设 $\xi_-(\cdot)$ 与 $\xi_+(\cdot)$ 分别表示与某个可容许解 $u$ 关联的、从任意点 $(\bar x,\bar t)\in(-\infty,\infty)\times(0,\infty)$ 出发的极小与极大后向特征线。则
\begin{align}
\begin{cases}
u(\xi_-(t)-,t)=u(\bar x-,\bar t)=u(\xi_-(t)+,t)\\
u(\xi_+(t)-,t)=u(\bar x+,\bar t)=u(\xi_+(t)+,t)
\end{cases}\quad 0<t<\bar t,\label{align:ch11-11-1-10}
\end{align}
\begin{align}
\begin{cases}
u_0(\xi_-(0)-)\leqslant u(\bar x-,\bar t)\leqslant u_0(\xi_-(0)+)\\
u_0(\xi_+(0)-)\leqslant u(\bar x+,\bar t)\leqslant u_0(\xi_+(0)+).
\end{cases}\label{align:ch11-11-1-11}
\end{align}
特别地, $u(\bar x+,\bar t)\leqslant u(\bar x-,\bar t)$ 对所有 $(\bar x,\bar t)\in(-\infty,\infty)\times(0,\infty)$ 成立, 且 $\xi_-(\cdot)$ 与 $\xi_+(\cdot)$ 重合当且仅当 $u(\bar x+,\bar t)=u(\bar x-,\bar t)$。
\end{theorem}

\begin{proof}
由\cref{chapter:ch10-广义特征线}可知, $\xi_-(\cdot)$ 与 $\xi_+(\cdot)$ 在 $[0,\bar t]$ 上均无激波。于是可应用\cref{theorem:ch11-11-1}, 取 $\sigma=0$, $\tau=\bar t$。由\eqref{align:ch11-11-1-4}, 当 $u(\bar x+,\bar t)=u(\bar x-,\bar t)$ 时, $\bar u=u(\bar x\pm,\bar t)$, 从而 $\xi_-(\cdot)$ 与 $\xi_+(\cdot)$ 重合。一般情形下, 考虑一个递增(或递减)序列 $\{x_n\}$, 收敛于 $\bar x$, 使得 $u(x_n+,\bar t)=u(x_n-,\bar t)$, $n=1,2,\cdots$。设 $\xi_n(\cdot)$ 表示从 $(x_n,\bar t)$ 出发的唯一后向特征线。则对所有 $t\in(0,\bar t)$, $u(\xi_n(t)\pm,t)=u(x_n\pm,\bar t)$。如\cref{chapter:ch10-广义特征线}所述, 序列 $\{\xi_n(\cdot)\}$ 从下方(或上方)收敛到 $\xi_-(\cdot)$(或 $\xi_+(\cdot)$)。因此, $u(\xi_-(t)-,t)=\lim u(x_n\pm,\bar t)=u(\bar x-,\bar t)$(或 $u(\xi_+(t)+,t)=\lim u(x_n,\bar t)=u(\bar x+,\bar t)$)。证毕。

\end{proof}

现在转向前向特征线的性质。

\begin{theorem}\label{theorem:ch11-11-4}
与可容许解 $u$ 关联的唯一前向广义特征线, 从任意点 $(\bar x,\bar t)\in(-\infty,\infty)\times(0,\infty)$ 出发。
\end{theorem}

\begin{proof}
假设两条不同的前向特征线 $\phi(\cdot)$ 与 $\psi(\cdot)$ 从 $(\bar x,\bar t)$ 出发, 使得对某个 $s>\bar t$, $\phi(s)<\psi(s)$。设 $\xi(\cdot)$ 表示从 $(\phi(s),s)$ 出发的极大后向特征线, $\zeta(\cdot)$ 表示从 $(\psi(s),s)$ 出发的极小后向特征线, 二者在 $[0,s]$ 上均无激波。对 $t\in[\bar t,s]$, $\xi(t)\geqslant\phi(t)$ 且 $\zeta(t)\leqslant\psi(t)$; 因而 $\xi(\cdot)$ 与 $\zeta(\cdot)$ 必在某个 $t\in[\bar t,s]$ 相交, 与\cref{corollary:ch11-11-2}矛盾。证毕。
\end{proof}

注意与之形成对照的是, 从 $x$ 轴上的点出发可以有多条前向特征线。特别地, 任何中心稀疏波的中心必定落在 $x$ 轴上。

下一个命题表明, 跳跃间断一旦形成, 便以激波的形式永久传播。

\begin{theorem}\label{theorem:ch11-11-5}
设 $\chi(\cdot)$ 表示与可容许解 $u$ 关联的唯一前向广义特征线, 从点 $(\bar x,\bar t)$ 出发, 其中 $\bar t>0$ 且 $u(\bar x+,\bar t)<u(\bar x-,\bar t)$。则对所有 $s\in[\bar t,\infty)$, $u(\chi(s)+,s)<u(\chi(s)-,s)$。
\end{theorem}

\begin{proof}
设 $\xi_-(\cdot)$ 与 $\xi_+(\cdot)$ 分别表示从 $(\bar x,\bar t)$ 出发的极小与极大后向特征线。由于 $u(\bar x+,\bar t)<u(\bar x-,\bar t)$, $\xi_-(\cdot)$ 与 $\xi_+(\cdot)$ 不同: $\xi_-(0)<\xi_+(0)$。

固定任意 $s\in[\bar t,\infty)$, 考虑从 $(\chi(s),s)$ 出发的极小与极大后向特征线 $\zeta_-(\cdot)$ 与 $\zeta_+(\cdot)$。对 $t\in[0,\bar t]$, 必有 $\zeta_-(t)\leqslant\xi_-(t)$ 且 $\zeta_+(t)\geqslant\xi_+(t)$。于是 $\zeta_-(0)<\zeta_+(0)$, 从而 $\zeta_-(\cdot)$ 与 $\zeta_+(\cdot)$ 不同。因此, $u(\chi(s)+,s)<u(\chi(s)-,s)$。证毕。
\end{proof}

由上, 可以确定激波发源的点。

\begin{definition}\label{definition:ch11-11-6}
称点 $(\bar x,\bar t)\in(-\infty,\infty)\times[0,\infty)$ 为\textbf{激波产生点}, 若从 $(\bar x,\bar t)$ 出发的某条前向广义特征线 $\chi(\cdot)$ 是激波, 即对所有 $t>\bar t$, $u(\chi(t)+,t)<u(\chi(t)-,t)$, 而从 $(\bar x,\bar t)$ 出发的每一条后向特征线均无激波。
\end{definition}

当 $(\bar x,\bar t)$ 是激波产生点且 $\bar t>0$ 时, 有两种可能: $u(\bar x+,\bar t)=u(\bar x-,\bar t)$ 或 $u(\bar x+,\bar t)<u(\bar x-,\bar t)$。前者中, 激波从 $(\bar x,\bar t)$ 出发时强度为零, 随演化而发展。后者中, 从 $(\bar x,\bar t)$ 出发有不同的极小与极大后向特征线 $\xi_-(\cdot)$ 与 $\xi_+(\cdot)$。夹在 $\xi_-(\cdot)$ 与 $\xi_+(\cdot)$ 图像之间的扇形区域必由特征线填满, 这些特征线将 $(\bar x,\bar t)$ 与 $x$ 轴连接, 由定义知它们无激波, 因而是直线。于是在这种情形下, 激波产生于中心压缩波的中心, 因而出发时即具有正强度。

\end{document}
